% Section draft for the paper (sec:canonicalquadrics), to be pasted into the Overleaf project.
% Companion of the pipeline stage FilterByCanonicalQuadrics (CanonicalQuadrics.m).
% Notation assumed from the paper: X = X_0^D(N)/W, the Atkin--Lehner group W, the genus g = g(X).

\section{Hyperellipticity from the canonical map}\label{sec:canonicalquadrics}

The tests of the preceding sections decide hyperellipticity from invariants that a hyperelliptic
curve cannot have: too many points over a finite field, a fixed-point count that no involution of a
hyperelliptic curve admits, a dual graph with no involution of tree quotient. None of them can
prove that a curve \emph{is} hyperelliptic, and for some quotients none of them applies. When the
holomorphic differentials of $X$ are available, the canonical map decides the question exactly in
both directions. This section records the criterion and the two situations in which we can compute
the differentials.

\subsection{Noether's count}

Let $C$ be a smooth projective curve of genus $g \ge 3$ over a field $k$ of characteristic $0$,
with canonical map $\phi\colon C \to \mathbf{P}^{g-1}$ given by a basis $\omega_1,\dots,\omega_g$ of
$H^0(C,\Omega^1)$. Max Noether's theorem states that if $C$ is not hyperelliptic then $\phi$ is an
embedding and the canonical ring is generated in degree one, so the multiplication map
\[
  \mu_2\colon \operatorname{Sym}^2 H^0(C,\Omega^1) \longrightarrow H^0(C,\Omega^{\otimes 2})
\]
is surjective; its target has dimension $3g-3$ by Riemann--Roch, its source dimension
$\binom{g+1}{2}$, so the kernel, the space of quadrics containing $\phi(C)$, has dimension
\[
  \binom{g+1}{2} - (3g-3) = \frac{(g-2)(g-3)}{2}.
\]
If $C$ is hyperelliptic, $\phi$ factors through the double cover $C \to \mathbf{P}^1$ followed by the
Veronese embedding of degree $g-1$, whose image is the rational normal curve; the quadrics through
it form a space of dimension $\binom{g-1}{2} = (g-1)(g-2)/2$, and they are exactly the kernel of
$\mu_2$ since the products $\omega_i\omega_j$ span the pull-back of
$H^0(\mathbf{P}^1,\mathcal{O}(2g-2))$, of dimension $2g-1$. The two counts differ by $g-2 \ge 1$.
Hence
\[
  \dim \ker \mu_2 =
  \begin{cases}
    (g-1)(g-2)/2 & \text{if $C$ is hyperelliptic,}\\
    (g-2)(g-3)/2 & \text{otherwise,}
  \end{cases}
\]
and computing the rank of the $\binom{g+1}{2}$ products $\omega_i\omega_j$ decides hyperellipticity.
The criterion is insensitive to the field of definition of the basis: we may compute over any field
containing $k$, and the rank is the same.

\subsection{Classical modular curves: $q$-expansions}

For $D = 1$ the holomorphic differentials of $X = X_0(N)/W$ are the weight-$2$ cusp forms on
$\Gamma_0(N)$ fixed by every Atkin--Lehner involution in $W$, through $f \mapsto f(q)\,dq/q$. The
product of two such forms is a weight-$4$ form on $\Gamma_0(N)$ (invariant under $W$, but we do
not need this), and a weight-$4$ form on $\Gamma_0(N)$ vanishing to order greater than the Sturm
bound $\tfrac{4}{12}[\mathrm{SL}_2(\mathbf{Z}) : \Gamma_0(N)]$ is zero. So a linear relation among
the products $f_if_j$ holds if and only if it holds between their $q$-expansions through that bound,
and the dimension of the kernel of $\mu_2$ is the number of products minus the rank of the matrix of
coefficients $(a_n(f_if_j))_{n \le B}$, $B$ any integer strictly larger than the bound. Everything
is computed exactly over $\mathbf{Q}$, from Magma's basis of $S_2(\Gamma_0(N))$ and the matrices of
the Atkin--Lehner operators.

The count reproduces Ogg's list \cite{Ogg74}: among the curves $X_0(N)$ of genus at least $3$,
exactly $N = 30, 33, 35, 39, 40, 41, 46, 47, 48, 59, 71$ are hyperelliptic. On the quotients left
open by the other tests it gives (list and table of the $D = 1$ verdicts); each of them agrees with
Magma's own computation on a model of the quotient (\texttt{X0NQuotient}, \texttt{IsHyperelliptic}).

\subsection{Shimura curves: the stored models}

For $D > 1$ there are no $q$-expansions, and we use the equations of
Section~\ref{sec:models}: a model of $X$ as a fibre product of double covers
$y_i^2 = f_i(t)$ of the star curve $X^*$, from which Magma computes a basis of
$H^0(X,\Omega^1)$ of the function field and the rank of the products. The verdict is only as good
as the model, so before the count the stage checks (i) that the factors of the fibre product are
consistent with a single coordinate $t$ (every triple $f_Uf_Vf_T$ a square in $\mathbf{Q}[t]$,
Section~\ref{sec:models}), (ii) that the function field has the genus the genus formula gives,
and (iii) that its numbers of points over $\mathbf{F}_p$ and $\mathbf{F}_{p^2}$ at three primes of
good reduction agree with the trace formula of Section~\ref{sec:trace}. A quotient stored directly as
$y^2 + h(t)y = f(t)$ is hyperelliptic by construction, and is accepted after the same genus and
point-count checks. At $N > 1$ the models rest on the value of the $m = 0$ term of the Borcherds
product (Section~\ref{sec:borcherds}), and the table marks those verdicts accordingly.

\subsection{Results}

On the (date) run the stage decides $35$ quotients of genus $3$ to $6$ that every other test leaves
open: $28$ at $D = 1$ by $q$-expansions ($8$ hyperelliptic, $20$ not) and $7$ at $D > 1$ from stored
fibre products (all of genus $3$, none hyperelliptic). (Table: $D$, $N$, $W$, $g$, verdict, number of
quadrics; the $D = 1$ rows checked against \texttt{X0NQuotient}, the $D > 1$ rows against the plane
models of the fibre products.) In particular the three quotients $X_0(114)/\langle w_2,w_{38}\rangle$,
$X_0(130)/\langle w_2,w_{26}\rangle$ and $X_0(195)/\langle w_5,w_{195}\rangle$ are hyperelliptic of
genus $3$, although earlier point-count certificates had suggested otherwise
(Section~\ref{sec:specialfiber}); the certificates were copied from a source at a different prime,
and the copying rule was corrected (Section~\ref{sec:specialfiber}).

\begin{remark}
Noether's theorem holds over any field of characteristic $0$, and the count makes no use of the
modular structure. The two obstacles to applying it everywhere are the differentials themselves,
which we only have at $D = 1$ or where a model exists, and, for the models, the independence of the
verdict from the model's correctness, which (i)--(iii) address but do not replace by a proof.
\end{remark}
